\documentclass[11pt]{article}
\usepackage[spanish]{babel}
\usepackage[utf8]{inputenc}
\usepackage[margin=2cm]{geometry}
\usepackage{amsmath}

\usepackage[american]{circuitikz}
\usetikzlibrary{calc}

\setlength{\parindent}{0pt}

\title{Boylestad -- Cap.~11: Problemas 14 y 17}
\date{}
\begin{document}
\maketitle

%==================================================
\section*{Problema 14 (Fig.~11.56): Amplificador de instrumentación}

\begin{center}
\begin{circuitikz}[scale=0.9, transform shape]
 % Op-amp superior (A1) y inferior (A2)
 \draw (0,4) node[op amp, noinv input up](A1){\texttt{A1}};
 \draw (0,0) node[op amp](A2){\texttt{A2}};
 % Entradas
 \draw (A1.+) -- ++(-1.2,0) node[left]{$V_1=3\,$V};
 \draw (A2.+) -- ++(-1.2,0) node[left]{$V_2=1\,$V};
 % Cadena resistiva
 \draw (A1.out) -- (2.5,4) coordinate(a);
 \draw (A2.out) -- (2.5,0) coordinate(b);
 \draw (a) to[R, l={$R=5\,\mathrm{k\Omega}$}, *-] (2.5,2.7) coordinate(c)
           to[vR, l={$R_P=1\,\mathrm{k\Omega}$}] (2.5,1.3) coordinate(d)
           to[R, l={$R=5\,\mathrm{k\Omega}$}, -*] (b);
 % Realimentación A1: inv -> c
 \draw (A1.-) -- ++(-0.4,0) |- (c);
 \draw (A2.-) -- ++(-0.4,0) |- (d);
 % Amp diferencial
 \draw (8,2) node[op amp](A3){\texttt{A3}};
 \draw (a) -- (4,4) to[R, l={$10\,\mathrm{k\Omega}$}] (5.5,4) |- (A3.-);
 \draw (b) -- (4,0) to[R, l_={$10\,\mathrm{k\Omega}$}] (5.5,0) -- (5.5,1.5) coordinate(p);
 \draw (p) -- (A3.+);
 \draw (p) to[short,*-] (6.3,1.5) to[R, l={$10\,\mathrm{k\Omega}$}] (6.3,-0.5) node[ground]{};
 \draw (A3.-) ++(-0.3,0) coordinate(n) -- ++(0,1.6) to[R, l={$10\,\mathrm{k\Omega}$}] (9.8,4.1) |- (A3.out);
 \draw (A3.out) to[short,-o] (10.5,2) node[right]{$V_o$};
\end{circuitikz}
\end{center}

\textbf{Fórmula (Boylestad, amp.\ de instrumentación):}
\[
|V_o| = \left(1+\frac{2R}{R_P}\right)(V_1-V_2)
\]

\textbf{Solución:}
\[
k = 1+\frac{2(5\,\mathrm{k\Omega})}{1\,\mathrm{k\Omega}} = 11
\qquad\Rightarrow\qquad
|V_o| = 11\,(3\,\mathrm{V}-1\,\mathrm{V}) = 22\ \mathrm{V}
\]

\textbf{Verificación nodo a nodo:}
\begin{align*}
I &= \frac{V_1-V_2}{R_P} = \frac{3-1}{1\,\mathrm{k\Omega}} = 2\ \mathrm{mA}\\
V_{o1} &= V_1 + I R = 3 + (2\,\mathrm{mA})(5\,\mathrm{k\Omega}) = 13\ \mathrm{V}\\
V_{o2} &= V_2 - I R = 1 - 10 = -9\ \mathrm{V}\\
V_o &= \frac{R_f}{R_1}\,(V_{o2}-V_{o1}) = 1\cdot(-9-13) = -22\ \mathrm{V}
\end{align*}

Como $V_{o1}$ entra por la entrada inversora de A3 (tal como está dibujada la figura), el signo es negativo:
\[
\boxed{V_o = -22\ \mathrm{V}\quad(|V_o| = 22\ \mathrm{V})}
\]

%==================================================
\newpage
\section*{Problema 17 (Fig.~11.59): Filtro pasabanda}

\begin{center}
\begin{circuitikz}[scale=0.9, transform shape]
 % Etapa 1: pasaaltas
 \draw (0,0) node[op amp, noinv input down](B1){};
 \draw (B1.+) -- ++(-0.6,0) coordinate(x1);
 \draw (x1) to[C, l_={$0.05\,\mu\mathrm{F}$}, -o] ++(-2.2,0) node[left]{$V_1$};
 \draw (x1) to[R, l={$10\,\mathrm{k\Omega}$}, *-] ++(0,-2) node[ground]{};
 \draw (B1.-) -- ++(-0.6,0) -- ++(0,1.6) coordinate(y1);
 \draw (y1) to[R, l_={$4.1\,\mathrm{k\Omega}$}, *-] ++(-2,0) node[ground]{};
 \draw (y1) to[R, l={$4.1\,\mathrm{k\Omega}$}] ++(3.3,0) coordinate(z1) |- (B1.out);
 % Etapa 2: pasabajas
 \draw (B1.out) to[short, *-] ++(0.3,0) to[R, l={$20\,\mathrm{k\Omega}$}] ++(2,0) coordinate(x2);
 \draw (x2) to[C, l={$0.02\,\mu\mathrm{F}$}, *-] ++(0,-2) node[ground]{};
 \draw (7,0.5) node[op amp, noinv input down](B2){};
 \draw (x2) -| ($(B2.+)+(-0.5,0)$) -- (B2.+);
 \draw (B2.-) -- ++(-0.5,0) -- ++(0,1.6) coordinate(y2);
 \draw (y2) to[R, l_={$20\,\mathrm{k\Omega}$}, *-] ++(-1.8,0) node[ground]{};
 \draw (y2) to[R, l={$20\,\mathrm{k\Omega}$}] ++(3.2,0) |- (B2.out);
 \draw (B2.out) to[short, *-o] ++(1,0) node[right]{$V_o$};
 \node at (0,-2.6) {Etapa 1: pasaaltas};
 \node at (6,-2.6) {Etapa 2: pasabajas};
\end{circuitikz}
\end{center}

\textbf{Fórmula (Boylestad, filtro RC de primer orden):}
\[
f = \frac{1}{2\pi R C}
\]

\textbf{Frecuencia de corte baja} (etapa pasaaltas: $R_1=10\,\mathrm{k\Omega}$, $C_1=0.05\,\mu\mathrm{F}$):
\[
f_{OL} = \frac{1}{2\pi (10\times10^{3})(0.05\times10^{-6})}
= \frac{1}{3.1416\times10^{-3}}
\;\Rightarrow\;
\boxed{f_{OL} \approx 318.3\ \mathrm{Hz}}
\]

\textbf{Frecuencia de corte alta} (etapa pasabajas: $R_2=20\,\mathrm{k\Omega}$, $C_2=0.02\,\mu\mathrm{F}$):
\[
f_{OH} = \frac{1}{2\pi (20\times10^{3})(0.02\times10^{-6})}
= \frac{1}{2.5133\times10^{-3}}
\;\Rightarrow\;
\boxed{f_{OH} \approx 397.9\ \mathrm{Hz}}
\]

\textbf{Ganancia en la banda de paso} (dato adicional):
\[
A_1 = 1+\frac{4.1\,\mathrm{k\Omega}}{4.1\,\mathrm{k\Omega}} = 2,\qquad
A_2 = 1+\frac{20\,\mathrm{k\Omega}}{20\,\mathrm{k\Omega}} = 2,\qquad
A_v = A_1A_2 = 4
\]

Banda de paso: $318.3\ \mathrm{Hz} \le f \le 397.9\ \mathrm{Hz}$ \quad (ancho de banda $\approx 79.6$ Hz).

\end{document}
